\chapter{文献与源码阅读方法}
\label{ch:reading}

\begin{keybox}
\textbf{本章把前面的知识转化成可复现的论文与源码阅读方法。}读论文时追踪“问题--算法--代价--证据”；读源码时追踪“入口--状态--控制流--数值操作--所有权变化--instrumentation”。任何源码结论都应绑定到具体 commit/tag，而不是只写“当前版本”。这样一年以后函数移动了，阅读笔记仍然可复现。
\end{keybox}

\section{论文阅读框架}
面对一篇 multigrid/HPC 论文，先写四行：求解什么 PDE/线性系统；采用什么 hierarchy、smoother 与 coarse solver；主要优化 compute、memory、sync 还是 communication；证据来自 convergence、time-to-solution、bandwidth 还是 scaling。四行写不出来，说明还没有抓住论文主问题。可以把这四行固定成一张阅读卡：
\begin{enumerate}
\item \textbf{数学对象：}PDE、离散、边界条件、矩阵性质是什么？
\item \textbf{多重网格对象：}hierarchy 怎样构造，smoother、$R/P$、coarse operator、bottom solver 分别是什么？
\item \textbf{执行对象：}数据放在哪里，哪个阶段发生同步/通信，粗层是否改变 rank/device 拓扑？
\item \textbf{证据：}作者证明或测量的是 convergence、cost per cycle、time-to-solution、bandwidth、strong/weak scaling 中的哪一个？
\end{enumerate}
如果一个结论无法落到这四格中的某一格，就先不要把它当成可迁移的工程结论。


\paragraph{先填一张假想论文的阅读卡。}
设摘要只告诉我们：“作者为三维扩散问题设计 GPU multigrid，并报告多 GPU scaling。”不要立刻追 kernel 名称。第一遍只填四格：数学对象是三维扩散离散；多重网格对象至少要继续查 smoother、transfer 和 coarse solver；执行对象是 GPU + 分布式 halo；证据对象应同时寻找 convergence 与 scaling。若论文正文后来发现两种实现使用不同 smoother 或容差，阅读卡就会立刻暴露“性能数字不可直接归因于硬件”这一问题。这个 worked example 的目的不是替你读论文，而是展示四格框架怎样主动生成后续问题。

\section{源码阅读卡：必须记录到函数与状态对象}

对具体软件实现，建议固定填写下面七格，而不是只画一个类图：

\begin{enumerate}
\item \textbf{源码快照：}repository、commit/tag、关键文件；
\item \textbf{public entry：}\texttt{solve/setup/apply} 从哪里进入；
\item \textbf{level state：}$A_\ell,x_\ell,b_\ell,r_\ell,P_\ell,R_\ell$ 分别存成什么对象；
\item \textbf{cycle owner：}谁真正决定 descend/bottom/ascend；
\item \textbf{operation implementation：}smooth、residual、restriction、prolongation 最终落到什么函数/kernel；
\item \textbf{ownership transition：}哪里发生 repartition、consolidation、device/host migration；
\item \textbf{instrumentation：}已有 profiler event/marker 在哪里。
\end{enumerate}

\paragraph{AMReX MLMG 的一张实际源码卡。}
以\chapref{ch:software}固定的 AMReX commit 为例：

\begin{itemize}
\item entry：\texttt{MLMGT<MF>::solve()}；
\item cycle control：\texttt{oneIter()} $\rightarrow$ \texttt{miniCycle()} $\rightarrow$ \texttt{mgVcycle()}；
\item numerical operations：\texttt{linop.smooth()}、\texttt{linop.restriction()}、\texttt{addInterpCorrection()}、\texttt{computeResOfCorrection()}；
\item bottom：\texttt{bottomSolve()} $\rightarrow$ \texttt{actualBottomSolve()}；
\item coarse ownership：\texttt{MLLinOpT::defineGrids()} 中的 agglomeration/consolidation 与 \texttt{BottomCommunicator()}；
\item profiling：\texttt{BL\_PROFILE("MLMG::oneIter()")}、\texttt{mgVcycle()}、\texttt{computeResidual()}、\texttt{actualBottomSolve()}。
\end{itemize}

这张卡填完以后，才适合继续追某个具体 smoother 或 GPU kernel。否则很容易在局部函数里看得很深，却不知道它位于整个 solve 的哪条依赖路径上。

\paragraph{hypre 与 PETSc 应怎样复用同一张卡。}
BoomerAMG 的 cycle owner 是 \texttt{hypre\_BoomerAMGCycle()}，level state 是 \texttt{A/P/R/U/F} arrays；PETSc PCMG 的 cycle owner 是 \texttt{PCMGMCycle\_Private()}，level state 则集中在 \texttt{PC\_MG\_Levels} 的 Mat/Vec/KSP 成员。对象名字不同，但七格问题完全相同。

\section{主题化阅读路线}

\subsection{多重网格数学基础}

建议先读 Brandt 的经典论文，理解多层方法为何有可能达到 $\Oh(N)$；然后用 Briggs--Henson--McCormick 的教程建立 two-grid/V-cycle 的可操作直觉；最后用 Trottenberg--Oosterlee--Schuller 深入 LFA、smoothing、coarse-grid approximation 与复杂问题\cite{brandt1977,briggs2000,trottenberg2001}。

阅读时不要先记 V-cycle 伪代码，而要反复回答：
\begin{equation}
\text{当前操作正在消除哪个误差子空间？}
\end{equation}

\subsection{块结构几何多重网格}

先用 AMReX 框架论文理解 MultiFab/BoxArray/DistributionMapping，再进入\chapref{ch:software}固定的源码快照\cite{zhang2019amrex,zhang2021amrex}。源码第一遍只追
\begin{equation}
\boxed{
\texttt{solve}
\rightarrow
\texttt{oneIter}
\rightarrow
\texttt{mgVcycle}
\rightarrow
\texttt{linop operations}
\rightarrow
\texttt{bottomSolve}.
}
\end{equation}
第二遍才进入 \texttt{MLLinOpT::defineGrids()}，观察 \texttt{BoxArray} 与 \texttt{DistributionMapping} 分别怎样被 agglomeration/consolidation 改写。官方 GPU/MLMG 文档适合补 API 与配置背景，但实现结论应以对应源码为准\cite{amrexGPUdocs2026,amrexMLMGdocs2026}。

\subsection{分布式代数多重网格}

BoomerAMG 的第一遍源码阅读直接从 \texttt{hypre\_BoomerAMGSetup()} 开始\cite{falgout2002hypre,hypreBoomer2026}，把真实调用链记录为
\begin{equation}
\texttt{CreateS}
\rightarrow
\texttt{Coarsen*}
\rightarrow
\texttt{Build*Interp}
\rightarrow
\texttt{BuildCoarseOperator/RAP}.
\end{equation}
第二遍进入 \texttt{hypre\_BoomerAMGCycle()}，确认 residual、restriction、prolongation 分别落到哪些 ParCSR Matvec/MatvecT，以及 \texttt{level/lev\_counter/cycle\_param} 怎样驱动状态机。只有做到这一步，才真正知道某个 AMG 参数改变的是 setup graph、transfer，还是 solve hot path。

\subsection{性能可移植框架}

阅读 MueLu/Kokkos 与 Ginkgo 时也不要停在“支持哪些 GPU”。对 MueLu，优先找 Hierarchy/Level/Factory：哪个 factory 构造 aggregate、prolongator 与 coarse matrix；对 Ginkgo，优先找 Executor/LinOp 与 multigrid level factory：同一个 operator 怎样绑定到不同 executor。文档用于找到入口，真正的实现比较仍应回到源码对象和调用链\cite{edwards2014kokkos,mueluDocs2026,ginkgo2022}。

\subsection{现代研究主题}

建议按问题而不是按年份读：
\begin{itemize}
\item \textbf{低精度层级：}mixed/three-precision AMG\cite{tsai2023mixed,tsai2023three,zong2024fp16}；
\item \textbf{新硬件单元：}Tensor Core AMG；论文读算法与评测，公开 AmgT/HYPRE fork 读 mBSR、SpGEMM 与 SpMV 的真实实现\cite{lu2024amgt,amgtRepo2026}；
\item \textbf{数据布局与通信协同：}fine-grain blocking GMG\cite{antepara2024brick}；
\item \textbf{无矩阵高算术强度：}GPU Stokes/high-order multigrid，以及 matrix-free algebraic $hp$-multigrid\cite{stokesgpu2025,ohm2026hpmg}；
\item \textbf{大规模粗层通信：}non-Galerkin coarse grids 与 process reduction\cite{falgout2014nongalerkin}。
\end{itemize}

\section{本章小结}
理论论文阅读应把结论拆成数学对象、multigrid 对象、执行对象与证据；源码阅读则必须进一步落到固定 commit、public entry、level state、cycle owner、operation implementation、ownership transition 与 profiler boundary。类名和 API 只是导航信息，真正可迁移的知识来自“谁拥有状态、谁驱动控制流、一次数学操作最终怎样执行”。

\section*{习题}

\subsection*{论文结构与证据拆解}
\begin{enumerate}[label=\arabic*.]
\item 任选一个开源 3D 7-point stencil kernel，从源码逐行统计每 cell 的显式 FLOP、数组读取/写入和理论 bytes，计算 arithmetic intensity，并与项目 benchmark 报告比较。
\item 从一个实际 V-cycle trace 中记录各 level unknowns 与时间，计算每 unknown 时间、相邻层工作量缩减比和 coarse overhead 比例。
\item 任选一篇论文中的 scaling 图，读取至少四个数据点，重新计算 speedup/efficiency；检查作者的结论是否与数据一致。
\item 对一个开源 AMG hierarchy 的日志，利用每层 unknowns 与 nnz 计算 grid/operator complexity，并与论文或文档报告对照。

\item 选择源码中的一个 residual/stencil kernel，把代码还原成数学算子，证明它与书中对应离散式等价（忽略边界处理差异）。
\item 选择一个 restriction/prolongation 实现，从索引和权重推导其数学矩阵 $R$ 或 $P$，并检查是否保持常数。
\item 选择一篇 multigrid 论文中的核心复杂度结论，依据论文给出的层级缩放假设重新完成推导，而不是直接引用作者结论。
\end{enumerate}

\subsection*{复现与阅读实践}
\begin{enumerate}[label=\arabic*.]
\item 任选一个开源 solver，成功构建并运行最小 Poisson case；记录 build 配置、问题规模、solver 参数和 residual history。
\item 任选 AMReX、hypre 或 PETSc，固定一个 commit/tag，按本章七格源码卡追踪一次 V-cycle：从 public entry 到 cycle driver，再到 pre-smooth、residual、restriction、coarse solve、prolongation、post-smooth，并标注 level state 与 ownership transition。
\item 在不改变数学算法的前提下加入一个计时器或 profiler marker，验证你能从运行结果定位到某一个书中定义的阶段。
\item 复现论文或官方文档中的一个小 benchmark；若无法复现绝对数字，至少比较趋势并分析硬件/版本差异。
\end{enumerate}

\subsection*{研究路线设计}
\begin{enumerate}[label=\arabic*.]
\item 用四格阅读卡整理一篇论文：“数学对象、多重网格对象、执行对象、证据”。任何无法落入四格的术语都必须继续追踪定义。
\item 若论文只报告 kernel bandwidth，没有 residual history、iteration count 或 time-to-solution，列出哪些性能结论可以成立、哪些不能成立。
\item 对一个开源 solver，回答 hierarchy 在哪里构造、transfer 在哪里执行、coarse topology 何时变化，并把每个答案交叉映射到本书相应章节。
\item 写一份 2 页源码审查报告：一个数值设计选择、一个性能设计选择、一个你会进一步验证的假设。所有结论必须附代码位置或实测证据。
\end{enumerate}
